\begin{center}6.\end{center}
\noindent\hspace*{\parindent}Die Begriffsbestimmung eines Integrals \(\int ft\,dt\) lässt sich auf zwei verschiedenen Arten feststellen.
\clearpage\SourceLeaf{420}\ArithmeticHead{413}{R}{CONVERGENZ DER ENTWICKLUNG PERIODISCHER FUNCTIONEN.}
\noindent\hspace*{\parindent}Einmahl erklärt man das Integral als eine solche Function von \(t\), aus deren Differentiation \(ft\,{.}\,dt\) hervorgeht.

Die zweite Art geht zunächst von einer Integration zwischen bestimmten Grenzen \(t_0\) und \(T\) aus. Man theilt die Differenz \(T-t_0\) in eine beliebige Anzahl, z. B. \(n\) gleicher oder ungleicher Theile
\[
t_1-t_0,\ t_2-t_1,\ t_3-t_2,\ { .}\,{.}\,{.},\ T-t_{n-1},
\]
\noindent und bestimmt das Aggregat
\[
\begin{gathered}
(t_1-t_0)ft_0+(t_2-t_1)ft_1+(t_3-t_2)ft_2\\
+\mathinner{\cdotp\,\cdotp\,\cdotp}+(T-t_{n-1})ft_{n-1}.
\end{gathered}
\]
\noindent Dieses Aggregat wird je nach der Art und Menge der Abtheilungen verschieden sein, aber sich einer bestimmten Grenze \(S\) unendlich annähern, wenn man die Grösse der Theile unendlich abnehmen, mithin ihre Anzahl unendlich zunehmen lässt. Mit dem Ausdruck Integral \(\int ft\,{.}\,dt\) von \(t=t_0\) bis \(t=T\) genommen, oder symbolisch mit
\[
\int\limits_{t_0}^{T} ft\,{.}\,dt
\]
\noindent bezeichnet man nun eben diesen Grenzwerth \(S\), welcher in der abgekürzten Sprache als Summe unendlich vieler unendlich kleiner Elemente \(ft\,{.}\,dt\) zwischen den Werthen \(t=t_0\) und \(t=T\) auftritt, oder in geometrischer Einkleidung als Quadratur. Da dieses \(S\) abhängig ist von der Wahl der Grenzwerthe \(t_0\) und \(T\) und bei gleich bleibendem \(t_0\) sich mit \(T\) zugleich verändert, so kann man es wie eine Function von \(T\) betrachten, und indem man diese Function mit \(FT\) bezeichnet, wird, wie sich leicht nachweisen lässt,
\[
dFT=fT\,{.}\,dT
\]
\noindent werden, oder indem man \(t\) anstatt \(T\) schreibt,
\[
dFt=ft\,{.}\,dt.
\]
\noindent Es ist hieraus klar, dass beide Begriffsbestimmungen zusammenfallen.
